\documentclass[12pt]{jlreq}
\usepackage{amsmath, amssymb}

\newcommand{\C}{\mathbb{C}}
\newcommand{\R}{\mathbb{R}}
\newcommand{\Z}{\mathbb{Z}}
\newcommand{\Tr}{\operatorname{Tr}}
\newcommand{\Impart}{\operatorname{Im}}
\newcommand{\AF}{\operatorname{AF}}
\newcommand{\EG}{\operatorname{EG}}
\newcommand{\AX}{\operatorname{AX}}
\newcommand{\GL}{\operatorname{GL}}
\newcommand{\Hom}{\operatorname{Hom}}
\newcommand{\Ker}{\operatorname{Ker}}
\newcommand{\Id}{\operatorname{Id}}
\newcommand{\Sym}{\operatorname{Sym}}

\begin{document}

$(R, \mathfrak{m})$ を可換なネーター局所環とし，$M, N$ を有限生成 $R$ 加群，$\phi: M \to N$ を $R$ 加群の準同型とする。また $r \in \mathfrak{m}$ は，任意の非零元 $x \in N$ に対し $rx \ne 0$ となるものとする。

\begin{enumerate}
  \item $K = \Ker\phi$ とおいたとき，自然な写像 $K/rK \to M/rM$ は単射になることを示せ。
  \item $\phi$ から誘導される写像 $\overline{\phi}: M/rM \to N/rN$ が同型写像ならば，$\phi$ も同型写像であることを示せ。
  \item $\mathbb{C}[[x, y, z]], \mathbb{C}[[t]]$ を $\mathbb{C}$ 上の形式的べき級数環とし，$f: \mathbb{C}[[x, y, z]] \to \mathbb{C}[[t]]$ を
  \[
    x \mapsto t^5,\quad y \mapsto t^6,\quad z \mapsto t^7
  \]
  で定義される $\mathbb{C}$ 上の環準同型とする。このときイデアル $\Ker f$ の有限個の元からなる生成系を一組求めよ。
\end{enumerate}

\end{document}